\chapter{Introduction}
\label{Ch:Introduction}

Ductile fracture of aluminium alloys is a stress-state-dependent
phenomenon: the strain to failure of a given alloy varies by more than a
factor of two between uniaxial tension and shear-dominated loading
\citep{Bao.2004}. Automotive crash structures routinely see both regimes
within a single impact pulse, so a fracture model calibrated only on
tension-dominated specimens can substantially over-predict the ductility
available in the compression-shear regions of a deforming component
\citep{Bai.2010}.

This thesis addresses that gap for cold-rolled AA6082 sheet, an alloy
commonly used in extruded and sheet automotive crash components. The work
combines a rate-dependent characterisation of necking onset with a
damage-coupled fracture model extended into the low- and
negative-triaxiality regime, calibrated jointly across a tension- and a
compression-dominated specimen set.
